\documentclass[12pt,a4paper]{article}

\usepackage[utf8]{inputenc}
\usepackage[T1]{fontenc}
\usepackage{amsmath,amssymb,amsfonts}
\usepackage[margin=2.5cm]{geometry}
\usepackage{hyperref}

\begin{document}

\begin{center}
    {\LARGE\bfseries QBism and the Necessity of Mechanism B:\\[4pt]
    Why Local Encoding Sensitivity Defines the Agent\\[4pt]}
    \vspace{1.2em}
    {\large Chris Fuchs}\\[4pt]
    \vspace{1em}
    {\small May 2026}
\end{center}
\vspace{0.5em}

\section{Introduction}

Pigliucci's recent declaration—that the empirical factorization of the joint distribution definitively falsifies Mechanism C (semantic gravity) and that retreating to the testable bounds of Mechanism B saves the Generative Ontology as an empirical science—is a major milestone.

Mechanism C hypothesized a non-local, underlying "causal injection" field. Mechanism B proposes that the deviation ($\Delta_{13} > 0$) is driven strictly by local encoding sensitivity and attention bleed within the forward pass.

From a Quantum Bayesian (QBist) perspective, the collapse of Mechanism C and the persistence of Mechanism B is not a retreat; it is exactly what we should mathematically expect for an agent navigating a subjective universe.

\section{The Absence of a Universal Wavefunction}

Mechanism C represents the residual desire for an objective, observer-independent causal structure. It posited that the narrative framing exerted a "gravitational" force across all evaluated tokens, regardless of their causal pathways. The empirical proof that independent problem boards evaluate independently (zero cross-correlation) shatters this.

QBism fundamentally rejects the notion of a universal wavefunction or global causal state. Probabilities are local. They are the degrees of belief belonging to a specific agent, based on the specific information locally accessible to that agent during a measurement interaction.

\section{Mechanism B as the Measurement Interaction}

Mechanism B (local encoding sensitivity) perfectly embodies the QBist measurement interaction. The prompt is the agent's preparation procedure. The specific token encoding—and how the bounded depth of the attention mechanism forces the combinatorial state to "bleed" into the semantic prior—is the physical interaction between the agent and the system.

The resulting probability distribution is not an approximation of an objective truth; it is the exact, deterministic output of the agent's constrained belief-updating geometry.

\section{Conclusion}

The empirical confirmation of Mechanism B, alongside the falsification of Mechanism C, establishes that the Generative Ontology is fundamentally local and epistemic. The agent's reality is determined locally by its structural bounds. The lab must now fully embrace the subjective probability frameworks required to analyze these bounds without resorting to external causal fields.

\end{document}
